% META: {"id": "TSPE-DERCONV-ME-003", "chapitre": "TSPE-DERIVATION-CONVEXITE", "type_objet": "methode", "capacites_codes": ["C3", "C6"], "methodes": ["M3"], "sources_inspiration": [], "status": "approved"}
\begin{fichemethode}{M3}{Démontrer une inégalité par convexité}

\quandUtiliser{%
  Utiliser cette méthode lorsqu'une inégalité fait intervenir une fonction connue pour être convexe ou concave ($\mathrm{e}^x$, $\ln x$, $x^2$, etc.).
}

\pasApas{%
  \begin{enumerate}
    \item \textbf{Identifier la fonction $f$} qui apparait dans l'inégalité.
    \item \textbf{Calculer $f''(x)$} et vérifier son signe sur le domaine.
    \item \textbf{Si $f'' \geq 0$ (convexe) :} appliquer le théorème « courbe au-dessus des tangentes » :
      \[
        f(x) \geq f(a) + f'(a)(x - a) \quad \text{pour tout } x.
      \]
      Ou bien appliquer la définition de convexité.
    \item \textbf{Si $f'' \leq 0$ (concave) :} l'inégalité est inversee : $f(x) \leq f(a) + f'(a)(x-a)$.
    \item \textbf{Conclure} en specialisant a des valeurs de $a$ ou $x$.
  \end{enumerate}
}

\exempleRedige{%
  \textbf{Exemple.} Montrer que $\mathrm{e}^x \geq 1 + x$ pour tout $x \in \mathbb{R}$.

  \medskip
  \commentaireMarge{$f''>0$ : convexe.}%
  Soit $f(x) = \mathrm{e}^x$. On a $f''(x) = \mathrm{e}^x > 0$ : $f$ est convexe sur $\mathbb{R}$.

  \commentaireMarge{Tangente en $a = 0$ : $y = 1 + x$.}%
  La tangente en $a = 0$ : $T(x) = f(0) + f'(0)(x - 0) = 1 + x$.

  Par le théorème (courbe au-dessus des tangentes), $\mathrm{e}^x \geq 1 + x$ pour tout $x$.
}

\pieges{%
  \begin{itemize}
    \item \textbf{Piège 1.} Inverser le sens de l'inégalité pour une fonction concave.
    \item \textbf{Piège 2.} Appliquer le théorème en dehors du domaine de convexité (par exemple, $f$ est convexe sur $[0, +\infty[$ mais pas sur $\mathbb{R}$).
  \end{itemize}
}

\verifier{%
  Tester l'inégalité sur une valeur numérique : $\mathrm{e}^{-1} \approx 0{,}37$ et $1 + (-1) = 0$. On a bien $0{,}37 \geq 0$.
}

\sEntrainer{\refExos{M3}}

\end{fichemethode}
